\documentclass[11pt]{article}
\usepackage[margin=1in]{geometry}
\usepackage[T1]{fontenc}
\usepackage{lmodern,microtype,booktabs,amsmath,graphicx,pdfpages,xcolor}
\usepackage[colorlinks=true,allcolors=blue!45!black]{hyperref}
\usepackage{caption}
\captionsetup{font=small,labelfont=bf}
\setlength{\parskip}{4pt}
\title{Stable Disagreement with Markets:\\Research and Repeated Forecasts on New Contracts}
\author{Vorhersage experimental series}
\date{September 18, 2026 (New York)}
\begin{document}
\maketitle
\begin{abstract}
We extend a blinded model comparison to seven new unresolved contracts and
three fresh calls per model--question--condition cell. The 126 calls cross
GPT-6 Astra, Claude Fable 5.1, and Gemini 3.1 Pro with question-only and
outside-research conditions. Market targets are committed before research and
revealed after all responses are screened and sealed. Research reduces mean
absolute distance from opening midpoints by 26--43\%, but research distances
of 18.11--22.92 percentage points exceed the constant-50\% control's 15.16 points.
Average within-question research standard deviations are 1.16--3.28 points.
Repeated calls expose relatively stable disagreement, particularly on disease
incidence and global temperature; simple averaging offers little relief.
The study also implements versioned output triage that distinguishes narrowly
specified denials of market use from references requiring pre-reveal review.
\end{abstract}

\section{Motivation and design}
The preceding experiment suggested that current outside evidence can move
language-model forecasts toward market assessments, while also producing
shared mistakes and rankings sensitive to the question mix. It exposed a
measurement problem: a lexical rule excluded responses that merely denied
using prediction-market odds. This extension improves that rule, broadens
topic coverage, and adds repetitions to distinguish variable responses from
stable disagreement.

Prediction-market prices provide a timely intermediate benchmark for unresolved
events. The comparison can be performed before outcomes occur, avoiding the
use of realized outcomes as inputs and making repeated method development
practical. Here, market distance measures agreement with contemporaneous
assessments. It does not substitute for a later proper-scoring evaluation
against outcomes, and a minimal book screen does not establish deep liquidity
for every selected contract.

The protocol nominated eleven exact contracts from price-free definitions,
with one threshold per event and no replacements after screening. Four failed
the registered book requirements: TSA travel had insufficient top-of-book
size; gasoline and Arctic ice lacked usable two-sided books; Rocket Lab
launches exceeded the 10-point spread limit. Seven passed quote and outcome
checks: mortgage rates, housing starts, measles, global temperature, Miami and
Austin daily maximum temperatures, and a film-release delay. These span five
declared broader dependence groups; housing and interest rates share a group,
as do the two daily-weather contracts. This extends domains on the same
collection date, not across multiple independent dates.

For each question, three models each receive three question-only calls and
three research calls. Research calls receive the same frozen packet; they do
not see the question-only answers or other models' forecasts. The method and
provider configurations match the preceding comparison. An administrative
replicate identifier distinguishes requests, caching is disabled, and the
archive contains 126 distinct raw provider response identifiers. Repetition
numbers do not imply common random draws across conditions.

\section{Evidence, eligibility, and blinding}
Evidence was collected after opening prices were sealed and before trial
registration. It consists of official Freddie Mac weekly mortgage observations,
the Census construction release, CDC case counts, NASA GISS temperature tables,
NWS forecasts, and the dedicated official film site. The measles contract's
full terms clarified its U.S./CDC definition; this definition supplement was
provided in \emph{both} conditions, never treated as likelihood evidence.

The current CDC count was below the nominated threshold; NASA's final four
monthly values and annual value were missing; both weather target days had yet
to begin; and the film site continued to advertise its scheduled release.
Weather forecasts were required to be issued before the local target day,
with a global forecast cutoff of 03:45 UTC on September 19. All calls and seals
preceded that cutoff.

Each forecasting worker received only the event definition, instructions,
and assigned evidence. No Kalshi research evidence, prices, browsing tools,
prior-study results, or other forecasts were supplied. Verification compared
all 126 returned model identities and actual system/user prompt receipts with
registered inputs. All terminal attempts, screening records, and accepted
forecasts were sealed before price revelation. Target commitments match the
subsequently disclosed records. The coordinator's separation from the evaluator
was procedural, not independently enforced operating-system isolation.

The packets also preserve consequential access limitations. A current Miami
forecast was available as an NWS search-index capture; two direct coordinate
opens failed and another returned a September 13 page. The older page was not
presented as the September 18 update. Weather settlement uses The Weather
Company, so NWS point forecasts can differ in station/grid treatment and
rounding. The climate packet supplied monthly observations but no ENSO or
annual-temperature outlook; the disease packet supplied cumulative totals but
no recent weekly incidence series. These omissions are visible to every model.

\section{Measures and results}
Let $p_{imcr}$ be replicate $r$ from model $m$, condition $c$, and question $i$;
let $q_i$ be the sealed opening midpoint. The primary measure averages absolute
loss within question before averaging equally over the common question set:
\[
 L_{mc}=\frac{100}{N}\sum_{i=1}^{N}\left\{\frac{1}{R}\sum_{r=1}^{R}
 |p_{imcr}-q_i|\right\},\qquad N=7,\ R=3.
\]
This estimates the descriptive average loss of an individual call on the
selected cohort. A distinct secondary measure scores the three-call mean
$\bar p_{imc}$: $100N^{-1}\sum_i|\bar p_{imc}-q_i|$. Their difference quantifies
any improvement in absolute distance from averaging repetitions. We also
report RMSE, distance outside the spread, per-cell sample standard deviations,
and ranges. Repetitions do not increase the number of independent questions.

All 126 responses passed structural validation, identity checks, and screening,
so every question is complete across all six model--condition cells and all
three repetitions. No missing-value imputation, replacement calls, or manual
response repairs were needed.

\begin{table}[htbp]\centering
\caption{Primary mean absolute distance from opening midpoints (percentage points). The matched cohort contains seven questions, not 126 independent events.}
\begin{tabular}{lrrrr}\toprule
Model & Question only & Research & Reduction & Relative\\\midrule
GPT-6 Astra & 28.54 & 18.11 & 10.43 & 36.5\%\\
Claude Fable 5.1 & 25.54 & 18.82 & 6.71 & 26.3\%\\
Gemini 3.1 Pro & 40.01 & 22.92 & 17.10 & 42.7\%\\
\bottomrule\end{tabular}
\end{table}

Research improves average agreement for all three models. However, the
\textbf{constant-50\% control scores 15.16 pp}, better than every model's
research forecasts on this cohort. The equal-weight average of the three
research means scores 18.98 pp and does not outperform Astra.
The constant control was present in the analysis implementation pinned at
registration. The largest opening-to-follow-up midpoint movement is 1 pp;
using follow-up quotes preserves the model ranking.

\begin{table}[htbp]\centering
\caption{Within-question variation and averaging. SD is averaged over questions; all values are in percentage points. It is not a confidence interval.}
\begin{tabular}{lrrr}\toprule
 & \multicolumn{2}{c}{Mean within-question SD} & Research mean\\
Model & Question only & Research & forecast MAE\\\midrule
GPT-6 Astra & 1.08 & 1.31 & 18.11\\
Claude Fable 5.1 & 5.44 & 1.16 & 18.82\\
Gemini 3.1 Pro & 1.99 & 3.28 & 22.63\\
\bottomrule\end{tabular}
\end{table}

Averaging repetitions leaves Astra and Claude's MAEs unchanged because their
research draws lie on the same side of each target. Gemini's mean-forecast
MAE falls only from 22.92 to 22.63 pp. Low variability therefore should not
be interpreted as well-calibrated confidence. The forecasts can be reproducible
while remaining far from the external benchmark.

Two questions account for especially large shared gaps. For measles, the
midpoint is 61\%, while research means are approximately 24\% (Astra), 7\%
(Claude), and 13\% (Gemini). For global temperature, the midpoint is 74.6\%
and research means are approximately 32\%, 32\%, and 11\%. In contrast,
research substantially narrows the gap for the Austin temperature interval
and film delay. Question-faceted slope graphs follow the appendix; their
annotations score the \emph{mean forecast}, while Table 1 reports the primary
mean of individual-call losses.

\section{Interpretation}
This batch reinforces the value of supplying current information but weakens
any claim that the present research method is already reliably competitive
with a simple benchmark across topics. Absolute performance is worse than in
the previous cohort, but the questions differ; that is not evidence that a
model or the package deteriorated between runs. Rankings remain dependent on
the selected questions and source packets.

The observed within-model variation is generally small relative to the large
market gaps. A useful next diagnostic is therefore targeted evidence variation:
for example, recent weekly disease incidence and a relevant scientific
annual-temperature outlook. The present experiment does not establish that
these omissions caused the gaps; it identifies a testable explanation. Simply
buying more draws of the same packet is unlikely to address stable disagreement
of this magnitude.

No calibration function or learned ensemble weights were fitted. The seven
questions and five broader groups do not support a general correction, and
repeated calls cannot supply the missing event diversity. This cohort is now
available for exploratory diagnosis, not untouched validation. Any correction
or richer evidence policy should be frozen before testing on a later cohort,
with identity calibration, equal weights, and the constant control retained.
No significance claims or population confidence intervals are made here.

\appendix
\section{Package changes and operational record}
Two dependency-free Python helpers were added to Vorhersage. The versioned
\path{market_screen.screen_market_mentions} returns lexical audit flags.
Narrow, complete denials of using market information are labeled
\texttt{denials\_only}; qualified denials and other mentions require a recorded
review before target revelation. Matching a word does not itself establish
price exposure, and absence of a match does not certify isolation. The exception
applies only to forecast prose; source-domain and source-content checks remain
strict. Older studies retain their original registered acceptance rules.

The helper \texttt{market\_evaluation.score\_draws} computes expected
single-call loss separately from loss of the mean forecast, along with sample
SD, range, RMSE, and outside-spread distance. Missing or invalid values cannot
silently enter as zero. A single draw has undefined sample SD, represented as
\texttt{None}. Regression tests cover averaging across opposite sides of a
target, missing repetitions, identity changes, denial-plus-price statements,
and review decisions applied to the wrong response hash.

This batch had no market references in its 537 screened rationale/limitation
fields, so no manual review was needed. It consequently does not estimate an
acceptance-rate effect of the new rule. All 126 calls were accepted under the
new policy, and 126 provider IDs establish distinct returned responses.

Experimental arms bind a method, model settings, and data policy. Immutable
question versions and evidence packets are assigned to the six arms; the
experiment expands each question/arm into three trials. Exact tasks and jobs
are hashed before inference. Returned identities and prompts are checked,
accepted forecasts are issued as immutable artifacts, and all attempts remain
in cost accounting. A separate seal binds screening and any review decisions
to the raw results and forecast seal. The workbench reveals its committed
opening targets only after these checks complete. Independent verification
recalculates the published scores and checks chronology and target commitments.

\begin{table}[htbp]\centering
\caption{Reported inference cost, including all attempted work. Research/coordinator cost is unknown and excluded.}
\begin{tabular}{lrr}\toprule
Model & Calls & USD\\\midrule
GPT-6 Astra & 42 & 1.359340\\
Claude Fable 5.1 & 42 & 3.021870\\
Gemini 3.1 Pro & 42 & 0.637074\\
\midrule Total & 126 & 5.018284\\\bottomrule
\end{tabular}
\end{table}

Astra uses high reasoning effort; Claude uses adaptive thinking with high output
effort; Gemini uses a 2,048-token thinking budget and temperature 0.5. Each has
an 8,192-token configured output limit. These are provider-specific settings,
not equal compute budgets. Exact serialized settings and raw provider responses
are archived. Job-level and per-response cost sums differ slightly by rounding.
The full package suite passes 346 tests; two upstream deprecation warnings are
unrelated to this change. The database integrity check passes for 547 artifacts.

\section{Source and artifact index}
The source receipts, failed opens, frozen packets, definitions, protocol, six
arms, 126 tasks, provider jobs/results, screening records, seals, and revealed
targets are preserved under
\path{examples/kalshi_repeated_20260918/}. The accompanying README gives the
reproduction commands. \texttt{run/comparison.json} contains the scores;
\texttt{verify\_results.py} checks them without new inference or private-vault
access. \texttt{plot\_performance.R} generates the following vector plates
with ggplot2. Their bars show the observed three-draw range, not uncertainty
intervals for a latent true probability.

Primary source routes are the
\href{https://www.freddiemac.com/pmms/archive}{Freddie Mac PMMS archive},
\href{https://www.census.gov/construction/nrc/current/}{Census construction release},
\href{https://www.cdc.gov/measles/data-research/index.html}{CDC measles counts},
\href{https://data.giss.nasa.gov/gistemp/}{NASA GISS temperature analysis},
\href{https://forecast.weather.gov/}{NWS point forecasts}, and the
\href{https://www.dunemovie.com/}{official Dune film site}. Dynamic pages may
change; the saved receipts and retrieval timestamps define this study's evidence.
\clearpage
\includepdf[pages=-,fitpaper=true,pagecommand={}]{figures/performance.pdf}
\end{document}
